In industrial monitoring, the cost of a sensor fault is paid not when it is detected but when it is acted upon. A reading that drifts away from its true value, a transmitter stuck at its last sample, or an instrument whose noise floor has grown all demand different maintenance responses: recalibration, replacement, cleaning, or process investigation. An operator who learns only that a signal is anomalous, and even which signal and in which direction, must still diagnose the failure mode before dispatching the correct action. The fault \emph{type} is what drives the maintenance action, and it is precisely what most streaming detectors leave to human triage.

The base detector we extend is a self-supervised, online anomaly detector for SCADA-based industrial systems~\citep{Wadinger2023}. It scores each signal by the conditional cumulative probability of its current value given the remaining signals, flags a value when that probability falls below $\alpha$ or above $1-\alpha$ (with $\alpha \approx 0.00133$, i.e.\ roughly three conditional standard deviations), and adapts its conditional model online with change-point awareness. Crucially, it is interpretable: it isolates the responsible signal and the direction of the excursion, and exposes per-signal operating limits. What it does not provide is the fault \emph{type}---it localizes the anomaly but does not characterize the failure mode, leaving the detection-to-action gap open.

We close that gap with a streaming sensor-fault classifier that attributes each flagged signal to one of four maintenance-relevant types---bias, drift, accuracy loss, or freezing---computed from the same conditional residuals the base detector already produces. The classifier maintains only a small set of exponentially weighted running statistics per signal: it uses O(1) memory, keeps no history buffers, and adds negligible cost to the streaming pipeline. The fault taxonomy itself is not new; we adopt the bias / noise / constant / drift categories established for sensor faults in wireless and environmental sensing~\citep{Sharma2010,Ni2009}. Our contribution is the \emph{streaming conditional-residual realization} of this taxonomy, composed directly with a self-supervised detector so that type attribution rides on the same online statistics as detection.

We lead with the practical outcome. On real battery-depletion faults from the Intel Berkeley Lab deployment~\citep{IntelLab2004}---where a mote's failing battery degrades its temperature and humidity readings to physically impossible values---the composed system reaches 90\% real-fault detection at a 7.4\% healthy false-positive rate, and types the resulting faults predominantly as freezing (0.93), the expected signature of a transmitter stuck at a degraded value. Online adaptation is not optional here: a static model fit on a short prefix of the same diurnal stream false-alarms at a rate of 0.97. The typing layer thus converts an adaptive detector's alarm into an actionable maintenance label on a real fault population.

We then characterize the operating envelope honestly, treating it as support for---not a qualification of---the practical result. Two claims we initially expected to hold do not survive our own validation, and we report them as refuted. First, freezing is \emph{not} a detection blind spot of the base detector: it catches every injected freeze on every channel, usually faster than a dedicated fixed-window freeze test, so the freeze test contributes type \emph{attribution} (a stuck sensor versus a process excursion), not detection capability. Second, the robustness of the adaptive detector to faults that arrive during adaptation comes from graded change-point suppression, not from conditional-residual cancellation. We treat this honesty as a feature: the envelope is mapped rather than overstated, and the limitations---bias absorption on strongly coupled signals, the limited cancellation of a two-signal conditional model, and the masking dead-band introduced by suppression---are quantified and explained.

The contributions of this paper are:
\begin{enumerate}
    \item A streaming, O(1)-memory sensor-fault classifier that types each flagged signal as normal, bias, drift, accuracy loss, or freezing from the conditional residuals of a self-supervised online detector, with no history buffers (Section~\ref{Method}).
    \item A demonstration on real battery-depletion faults that the composed system reaches 90\% real-fault detection at a 7.4\% healthy false-positive rate and assigns maintenance-relevant types, with online adaptation shown to be mandatory (Section~\ref{Results}).
    \item A rigorous characterization of the operating envelope on controlled injections, including the monotone dependence of bias recall on conditional coupling, and two refuted claims (freezing as a blind spot; residual cancellation as the source of adaptation robustness) reported as such (Sections~\ref{Results} and~\ref{Discussion}).
\end{enumerate}
